% La residencia capitular pertenece a la estructura maestra. Este archivo
% desarrolla su contenido y conserva una segunda etiqueta compatible.
\label{chap:realizacion-espectral-numero-particula}
\index{número!realización espectral}
\index{partícula!sección persistente}
\index{ruta!sombra observable}
\index{espectro!realización física}

El cierre de HMT produce números como secciones genealógicas completas. La
Mecánica Dimensional comienza cuando una sección de esa clase se realiza sobre
una ruta, transporta un registro, induce un operador espectral y adquiere una
carta física. Este paso conserva el antecedente matemático y cambia el
codominio: número, estado, ruta, espectro y partícula son objetos enlazados
por mapas explícitos y conservan tipos diferentes.

La bisagra formal se expresa mediante
\begin{equation}
 \boldsymbol x
 \xmapsto{\operatorname{sh}}
 \gamma_{\boldsymbol x}
 \xmapsto{\mathcal L}
 \mathcal L_{\boldsymbol x}
 \xmapsto{\Sigma}
 \sigma_{\boldsymbol x}
 \xmapsto{\chi}
 \chi_{\boldsymbol x}
 \xmapsto{\mathcal T}
 \widehat K_{\boldsymbol x}
 \xmapsto{\mathfrak R_{\rm phys}}
 \mathfrak p_{\boldsymbol x}.
 \label{eq:bisagra-numero-espectro-particula}
\end{equation}
Cada flecha posee un dominio propio y conserva la clave de procedencia de la
sección inicial.

\section{Correspondencia física del núcleo APP--TRIT--TPK}
\label{sec:reaplicacion-fisica-app-trit-tpk}

La apertura de Mecánica Dimensional vuelve a recorrer el núcleo formal con
un nuevo codominio. Esta repetición conserva los teoremas de HMT y declara
los mapas de realización física que se añaden. Su orden causal es
\[
 \boxed{\begin{gathered}
 \text{APP como soporte de hojas}
 \longrightarrow \text{TRIT como orientación temporal}\\
 \downarrow\\[-1mm]
 \text{ciclo unitario y defecto torsional}\\
 \downarrow\\[-1mm]
 \text{coordenadas angulares duales}
 \longrightarrow \text{partícula--ruta}
 \end{gathered}}
\]

\subsection{El toro de Cayley como soporte de hojas físicas}

APP aporta el toro bidimensional
\[
 X_9=(\mathbb Z/9\mathbb Z)^2
\]
con su grafo de Cayley y dos evaluaciones sobre los mismos vértices. En el
nivel formal, \(\Sigma\) acumula y \(\Pi\) pliega. La realización física
asigna a esas evaluaciones hojas y canales, conservando el mapa de
procedencia que permite comparar sus respuestas. El toro mantiene la
periodicidad espacial; la categoría de caminos mantiene dirección,
reversibilidad y memoria de transporte.

\subsection{Carta gnomónica euclídea y atlas de curvaturas}
\label{sec:carta-gnomonica-equivalencia-ads-ds}

El gnomón en L completa el núcleo multiplicativo \(8\times8\) y separa el
interior plegado de la carta de publicación. En
la clase nula se distinguen precisamente los dos lectores que el cociente
identifica fuera de la frontera,
\[
 n\longmapsto n\bmod 9,
 \qquad
 n\longmapsto\operatorname{dr}_9(n),
\]
de modo que el residuo \(0\) y su publicación nonádica \(9\) conservan
coordenadas distintas del mismo cierre. El sector \(3,6,9\) hace visible
esta transición entre la raíz digital y el residuo modular.

En la realización geométrica, la L gnomónica porta la carta euclídea
\(\mathscr H_0^{(L)}\). Sobre ella, los lectores inercial y gravitatorio
descienden a una normalización común,
\[
 \frac{\mathcal G_{L}}{\mathcal I_{L}}=1,
\]
que constituye la residencia del principio local de equivalencia. Antes de
esa proyección, el estado enriquecido conserva separadamente las lecturas
areal y volumétrica. El principio de Ambivalencia se expresa entonces por
la razón interna
\[
 \boxed{
 \frac{\mathcal L_{\rm ar}(x)}{\mathcal L_{\rm vol}(x)}
 =\frac{\pi_{\rm HMT}}{\varphi_{\rm HMT}^{2}}.}
\]
La igualdad euclídea y la razón ambivalente pertenecen a niveles distintos:
la primera es la proyección local sobre la L; la segunda conserva la
genealogía areal--volumétrica anterior a esa proyección.

La realización métrica se obtiene mediante una familia declarada de mapas.
Para
\[
 \mathbb A_s:=\mathbb R[X]/(X^2+s),
 \qquad s\in\{-1,0,+1\},
\]
y una escala de curvatura \(\rho>0\), se fija
\[
 \mathcal M_{s,\rho}:\mathbb A_s\dashrightarrow
 (M_s,g_{s,\rho}),
 \qquad
 g_{s,\rho}=\mathrm d r^2+S_{s,\rho}(r)^2\mathrm d\theta^2,
\]
con
\[
 S_{s,\rho}(r)=
 \begin{cases}
  \rho^{-1}\sinh(\rho r),&s=-1,\\
  r,&s=0,\\
  \rho^{-1}\sin(\rho r),&s=+1.
 \end{cases}
 \qquad
 K(M_s,g_{s,\rho})=s\rho^2.
\]
El dominio radial es \(r\geq0\) para \(s=-1,0\); para \(s=+1\), la carta
polar cubre \(0\leq r\leq\pi/\rho\), con la regularidad polar usual en sus
extremos. La rama \(s=-1\) suministra la sección espacial hiperbólica de la
carta de tipo anti--de Sitter; \(s=0\) realiza la carta euclídea gnomónica;
y \(s=+1\) proporciona la geometría espacial positiva que se prolonga en
las foliaciones de de Sitter. La correspondencia completa se recoge en los
\cref{sec:ansatz-metrico-tres-curvaturas,sec:atlas-ads-euclideo-ds-hmt}.

La clasificación algebraica y la familia métrica quedan coordinadas mediante
la realización declarada \(\mathcal M_{s,\rho}\). Esta flecha constituye el
ansatz de interfaz física, mientras
\(K=s\rho^2\) es el cálculo exacto de la métrica una vez fijados \(s\) y
\(\rho\). La relación cuadrática fija la clase algebraica; la realización
declarada transporta su signo a la curvatura, y \(\rho\) fija la escala
geométrica. En la carta negativa, la restricción al borde coordina la
proyección arquimediana con la incidencia excepcional; en la carta nula, los
lectores inercial y gravitatorio descienden a la equivalencia local; en la
carta positiva, las foliaciones satisfacen \(\ddot a/a=H^2>0\). Este atlas
enlaza la holografía de curvatura negativa, la física local euclídea y la
evolución cosmológica acelerada mediante dominios y mapas de cambio
declarados.

\subsection{Representación temporal del TRIT: (+1) retorno y memoria, (0) presente de umbral y (-1) apertura bidireccional}

Una vez declarado un orden sobre las transiciones, el lector temporal es
\[
\begin{aligned}
 +1&\longmapsto\text{retorno, memoria retenida y pasado},\\
 0&\longmapsto\text{umbral de evaluación y presente},\\
 -1&\longmapsto\text{apertura bidireccional y futuro}.
\end{aligned}
\]
El pasado designa el cierre transportado; el presente, la frontera de
evaluación; el futuro, el árbol de prolongaciones compatibles. Esta lectura
es la realización temporal de las tres álgebras cuadráticas ya construidas
en HMT y conserva la involución que intercambia apertura y retorno.

\subsection{Torsión mínima y dinámica unitaria en \texorpdfstring{\(\mathbb C^{108}\)}{C108}}

El invariante denominado torsión mínima es
\[
 \Delta_4
 =
 \frac{\pi-e\log\pi}{270}
 \left(1+\frac{\pi}{729}\right).
\]
Sus entradas proceden de las secciones HMT de \(\pi\) y \(e\), del canal
\(270\) y del interior \(729\). Funciona como memoria global del
refinamiento antes de la especialización electrónica.

La dinámica temporal de referencia reside en
\(\mathcal H_{\rm clk}=\mathbb C^{108}\). Con
\[
 \omega_0=\frac{2\pi}{108\,t_0},
 \qquad
 H_{\rm clk}
 =\hbar_{\rm int}\omega_0
 \sum_{k=0}^{107}k\,|\widetilde k\rangle
 \langle\widetilde k|,
\]
el propagador satisface
\[
 U_{\rm step}
 =\exp\!\left(-\frac{\mathrm i}{\hbar_{\rm int}}
 H_{\rm clk}t_0\right),
 \qquad
 U_{\rm step}^{108}=I.
\]
El ciclo unitario es exacto para las primitivas declaradas. La realización
como cristal de tiempo físico añade el Hamiltoniano o protocolo impulsado,
la respuesta subarmónica primitiva, la estabilidad bajo perturbaciones, la
persistencia y el mecanismo de protección desarrollados en el
\cref{chap:regimenes-cuadraticos-tiempo}. Esta residencia conserva la
jerarquía y remite allí la demostración.

\subsection{Coordenadas angulares duales y canales espectrales}

La incompatibilidad suma--producto se transporta a dos coordenadas angulares
\(A\) y \(C^*\). En la carta arquimediana
\[
 x=\frac{\pi A}{180},
 \qquad
 y=\frac{\pi C^*}{180},
\]
los canales son los autovalores del elemento central
\[
 T=e^{-xI-yR}=q_+P_++q_-P_-,
 \qquad
 q_-=e^{-x+y},
 \qquad
 q_+=e^{-x-y}.
\]
La carta es invertible:
\[
 x=-\frac12\log(q_+q_-),
 \qquad
 y=\frac12\log\frac{q_-}{q_+}.
\]
Por tanto, \(A\) y \(C^*\) son, respectivamente, las coordenadas simétrica y
antisimétrica de una única banda espectral. El
\cref{chap:canales} contiene su construcción y sus pruebas; la
presente bisagra fija su posición causal antes de la partícula--ruta.

\section{Estado, ruta, forma espectral y realización física}

\begin{definition}[Estado finito y sección genealógica]
\label{def:estado-finito-seccion-genealogica-md}
Un estado finito es un elemento
\(x_n\in\mathcal S_n^{\rm surv}\). Una sección genealógica es una familia
\(\boldsymbol x=(x_n)_n\) compatible a toda profundidad. El estado es una
restricción de la sección; la sección determina su sistema completo de
restricciones.
\end{definition}

\begin{definition}[Ruta observable]
\label{def:ruta-observable-bisagra-md}
Una carta secuencial de profundidad \(R\) define la sombra
\[
 \gamma_{\boldsymbol x}^{[R]}
 :=\operatorname{sh}_R(\boldsymbol x)
 =(c_0\to c_1\to\cdots\to c_R).
\]
La ruta conserva posición, hoja, orientación, acarreo y memoria según la
lista de datos de la carta. Es una presentación ordenada de la sección.
\end{definition}

\begin{definition}[Forma espectral]
\label{def:forma-espectral-numero-md}
La forma espectral asociada a una sección es
\[
 \mathfrak S(\boldsymbol x)
 :=(\gamma_{\boldsymbol x},
    \mathcal L_{\boldsymbol x},
    \sigma_{\boldsymbol x},
    \chi_{\boldsymbol x},
    \widehat K_{\boldsymbol x}),
\]
donde \(\mathcal L_{\boldsymbol x}\) es el registro de ruta,
\(\sigma_{\boldsymbol x}\) su firma entera, \(\chi_{\boldsymbol x}\) el
carácter construido sobre la firma y \(\widehat K_{\boldsymbol x}\) el
operador obtenido por la jerarquía de torres. El espectro
\(\operatorname{Spec}\widehat K_{\boldsymbol x}\) es una salida interna.
\end{definition}

\begin{definition}[Realización física]
\label{def:realizacion-fisica-forma-espectral}
Una realización física de \(\mathfrak S(\boldsymbol x)\) consiste en
\[
 \mathfrak R_{\rm phys}
 =(\mathcal H_{\rm phys},\rho_{\rm sym},\varsigma_{\bf u},\mathcal O),
\]
donde \(\mathcal H_{\rm phys}\) es el espacio de estados,
\(\rho_{\rm sym}\) la representación del grupo de simetría,
\(\varsigma_{\bf u}\) la carta de unidades y \(\mathcal O\) el conjunto de
observables. La realización entrelaza los transportes de ruta con la dinámica
de \(\mathcal H_{\rm phys}\) y publica magnitudes con unidades.
\end{definition}

La forma espectral es exacta dentro del sistema discreto una vez fijados sus
lectores. La identificación con una especie física pertenece al mapa de
realización. Esta separación permite publicar el operador, su espectro, su
posición y su estatuto, acompañados del entrelazador físico correspondiente.

\section{\(4\) gramáticas y sus dominios}

El paso desde el número hasta la partícula utiliza cuatro gramáticas. Las
tres primeras pertenecen al núcleo matemático; la cuarta pertenece a la
representación estadística elegida.

\begin{center}
\begin{tabular}{@{}lll@{}}
\toprule
\textbf{Gramática} & \textbf{Operación} & \textbf{Función}\\
\midrule
acumulación & hoja suma de APP & compone contribuciones aditivas;\\
pliegue & hoja producto de APP & compone escalas y transportes;\\
orientación & TRIT y conjugación de ruta & distingue sentido y anti-sección;\\
ocupación & proyector de la representación & fija multiplicidad estadística.\\
\bottomrule
\end{tabular}
\end{center}

La ocupación fermiónica se realiza, por ejemplo, en un álgebra CAR fijada:
\[
 \{a_i,a_j\}=0,
 \qquad
 \{a_i^\dagger,a_j^\dagger\}=0,
 \qquad
 \{a_i,a_j^\dagger\}=\delta_{ij}I.
\]
De estas premisas se deduce
\[
 (a_i^\dagger)^2=0,
 \qquad
 n_i:=a_i^\dagger a_i,
 \qquad
 n_i^2=n_i,
 \qquad
 \operatorname{Spec}(n_i)=\{0,1\}.
\]
El TPK selecciona la graduación y transporta la ruta; la representación CAR
realiza la exclusión. Esta distinción fija positivamente la procedencia
algebraica del principio de Pauli.

\section{El functor de realización espectral}

Sea \(\Omega_{\rm HMT}^{\rm surv}\) la categoría cuyos objetos son secciones
supervivientes y cuyos morfismos son transformaciones compatibles con las
restricciones. Sea \(\mathbf{SpecRoute}\) la categoría de paquetes
\((\gamma,\mathcal L,\sigma,\chi,\widehat K)\) y entrelazadores que
conservan las claves de ruta.

\begin{theorem}[Realización espectral de una sección]
\label{thm:functor-realizacion-espectral-numero}
Los lectores de ruta, registro, firma, carácter y torres definen, en su
dominio común, un functor
\[
 \mathfrak S:
 \Omega_{\rm HMT}^{\rm surv}\longrightarrow\mathbf{SpecRoute}.
\]
Para una restricción compatible de secciones, la ruta se trunca, el registro
se restringe y los lectores posteriores conmutan con esa restricción. Dos
presentaciones relacionadas por una simetría genealógica producen formas
espectrales entrelazadas.
\end{theorem}

\begin{proof}
La sombra de ruta está definida por lectura secuencial del estado y conmuta
con la truncación. El registro concatena datos locales tipados y su
restricción elimina únicamente el segmento posterior. La firma y el carácter
se calculan mediante homomorfismos declarados sobre ese registro. La
jerarquía de torres prolonga el carácter por composición. Identidad y
composición se preservan en cada etapa.
\end{proof}

El teorema precisa la expresión «el número adopta forma espectral»:
transporta una sección completa a un paquete de ruta y espectro que conserva
su genealogía.

\section{Persistencia y definición de partícula}

Para un prefijo \(s_N\in\mathcal S_N^{\rm surv}\), sea
\[
 \operatorname{Ext}_M(s_N)
 :=\{s_M\in\mathcal S_M^{\rm surv}:\rho_{M,N}(s_M)=s_N\},
 \qquad
 \mathcal L(s_N)
 :=\sup\{M:\operatorname{Ext}_M(s_N)\neq\varnothing\}.
\]

\begin{definition}[Persistencia]
\label{def:persistencia-bisagra-numero-particula}
Un prefijo es persistente cuando pertenece a una sección global compatible;
en un árbol finitamente ramificado de prolongaciones, esta condición equivale
a \(\mathcal L(s_N)=\infty\).
\end{definition}

\begin{definition}[Partícula matemática]
\label{def:particula-matematica-bisagra}
Sea \(\gamma_{\boldsymbol x}\) una ruta persistente y sea
\(\mathcal B_{\gamma_{\boldsymbol x}}\) un fibrado discreto cuyas fibras
portan los grados internos de la realización. Una partícula matemática es
\[
 \mathfrak p=
 (\boldsymbol x,
  \gamma_{\boldsymbol x},
  \mathcal B_{\gamma_{\boldsymbol x}},
  s,n,U,
  \mathfrak S(\boldsymbol x),
  \mathfrak R_{\rm phys}),
\]
donde \(s\) es una sección compatible del fibrado, \(n\) su función de
ocupación, \(U\) el propagador que entrelaza los transportes y
\(\mathfrak R_{\rm phys}\) la realización física declarada.
\end{definition}

La definición contiene tres cierres:
\begin{enumerate}
\item el número es el cierre global de prefijos compatibles;
\item la forma espectral es el cierre de la lectura de ruta bajo registro,
      firma, carácter y torres;
\item la partícula es el cierre local recorrido por una sección persistente
      de un fibrado físico.
\end{enumerate}
La relación se resume como
\[
 \boxed{
 \text{número}=\text{cierre genealógico global},
 \qquad
 \text{partícula}=\text{realización local persistente de ese cierre}.}
\]
La igualdad verbal designa una dependencia functorial entre tipos.

\section{Banda orientada y anti-sección}

Sea \(\tau=(\tau_0,\ldots,\tau_{11})\in\{-1,0,+1\}^{12}\) la palabra
orientada de una ruta y defínase
\begin{equation}
 C_M(\tau):=-\operatorname{rev}(\tau).
 \label{eq:conjugacion-material-bisagra}
\end{equation}
Se cumple \(C_M^2=I\). La banda
\[
 \mathcal B(\tau):=\{\tau,C_M\tau\}
\]
reúne dos orientaciones de una misma genealogía. Una sección realiza la
partícula y la sección conjugada realiza la antipartícula.

\begin{proposition}[Paridad de masa y orientación]
\label{prop:paridad-masa-orientacion-bisagra}
Sea \(\mathfrak m(\tau)\) una lectura de masa que dependa de invariantes pares
de la banda, y sea \(q(\tau)\) una lectura orientada impar. Entonces
\[
 \mathfrak m(C_M\tau)=\mathfrak m(\tau),
 \qquad
 q(C_M\tau)=-q(\tau).
\]
\end{proposition}

\begin{proof}
La primera identidad expresa la invariancia bajo la involución de banda. La
segunda expresa la paridad de una forma orientada. La relación \(C_M^2=I\)
garantiza la compatibilidad de ambas lecturas.
\end{proof}

Partícula y antipartícula comparten así el espectro de masa y poseen
orientaciones opuestas. La conclusión procede de la paridad de los lectores.

\section{Propagación y suma finita de caminos}

Sea \(U\) un propagador sobre un espacio finito de estados. Para todo
\(R\geq1\), la multiplicación matricial da
\begin{equation}
 (U^R)_{fi}
 =\sum_{\substack{\gamma:i\to f\\|\gamma|=R}}
   \prod_{\ell\in\gamma}U_\ell.
 \label{eq:feynman-finito-bisagra}
\end{equation}

\begin{theorem}[Expansión exacta por historias]
\label{thm:feynman-finito-bisagra}
La ecuación \eqref{eq:feynman-finito-bisagra} enumera una vez cada ruta de
longitud \(R\) entre \(i\) y \(f\). Si \(U\) es unitario, \(U^R\) conserva
la norma y define una evolución finita unitaria.
\end{theorem}

\begin{proof}
Al expandir el producto matricial aparecen todos los índices intermedios
\((x_1,\ldots,x_{R-1})\). Cada tupla determina una ruta y cada ruta determina
una tupla. La unitariedad se sigue de
\((U^R)^\dagger U^R=I\).
\end{proof}

La amplitud de onda es la suma coherente sobre la banda de rutas; la
partícula es la sección persistente cuya propagación realiza esa banda. Onda
y estado forman niveles complementarios de la misma realización.

\section{Horizonte finito de prolongación y clase de extensiones compatibles del estado genealógico}

\begin{definition}[Ruta virtual finita]
\label{def:ruta-virtual-finita-bisagra}
Un prefijo \(s_N\) es virtual finito si existe \(L<\infty\) tal que
\[
 \operatorname{Ext}_L(s_N)\neq\varnothing,
 \qquad
 \operatorname{Ext}_{L+1}(s_N)=\varnothing.
\]
Su ruta transporta registro, firma y amplitud hasta \(L\), y su horizonte de
prolongación termina en la frontera \(L+1\).
\end{definition}

En un árbol finitamente ramificado, la existencia de prolongaciones a toda
profundidad equivale, por el lema de König, a la existencia de una sección
global. La persistencia y la virtualidad quedan definidas mediante la
geometría de prolongaciones. La ruta virtual ocupa un estado intermedio
tipado y posee una duración genealógica finita.

\section{Espectro interno, realización y especie}

Para una multisección \(F\) de rutas, el objeto natural es el operador
\[
 \widehat B_F e_\Gamma
 =\kappa(\Gamma)e_\Gamma,
 \qquad
 \exp((\log R_{\rm act})\widehat B_F).
\]
Su espectro es invariante bajo permutaciones de la base de la fibra. La carta
física determina qué combinación de autovalores corresponde a un observable
escalar o matricial. Deben conservarse tres niveles:
\[
 \operatorname{Spec}\widehat B_F,
 \qquad
 \mathfrak R_{\rm phys}(\widehat B_F),
 \qquad
 \text{identificación de la especie}.
\]
El primero es una salida interna; el segundo es una realización; el tercero
es una clave dentro del inventario físico. Una tabla científica completa
publica los tres y el morfismo que los relaciona.

\section{Functor de realización físico-espectral y 1.ª clausura electrónica}

La bisagra queda fijada por la cadena
\[
 \boxed{\begin{gathered}
 \text{número genealógico}
 \longrightarrow\text{ruta observable}
 \longrightarrow\text{registro}
 \longrightarrow\text{firma}
 \\
 \text{carácter}
 \longrightarrow\text{torres}
 \longrightarrow\text{operador espectral}
 \\
 \text{realización física}
 \longrightarrow\text{partícula}
 \end{gathered}}
\]
La primera clausura física especializa esta cadena en la sección electrónica
seleccionada por el centro del atlas, con la recta de acción, los lectores
bidireccionales y la carta dimensional vigentes. Esta residencia establece el
dominio formal que permite comprender por qué el electrón realiza la teoría
del número y por qué su generalización a las demás especies conserva ruta,
espectro y genealogía.

\paragraph{Composición de los lectores dimensionales y compatibilidad entre sus invariantes matemáticos}
La definición de sección persistente y fibrado discreto pertenece al capítulo
de partículas; la ruta, el registro, la firma y el carácter, al capítulo de
extensión superviviente; la banda, la conjugación material y la virtualidad,
al capítulo de anti-secciones; los operadores de fibra y sus espectros, a la
ley bidireccional de masas. Las cuatro gramáticas y el cierre local del
recorrido se integran aquí con la estructura CAR, la persistencia, el carácter
y la realización física.
